\section{Width compression with controlled product norm}
\label{sec:pf-width-compression}

The width dependence in Theorem~\ref{thm:product-target-log-padding}
can be removed. Boundedness and near-linearity of tanh permit a
norm-preserving compression: very large normalized rows can be dropped,
very small rows can be linearized, and the remaining neurons can be
sampled with a probability chosen to balance their two adjacent norms.

\begin{theorem}[Finite-set width compression for tanh]
\label{thm:pf-width-compression}
Let $X$ be a nonempty finite set of inputs of norm at most $R\ge1$.
Suppose a depth-$D_0\ge2$ tanh network with linear readout, arbitrary
finite hidden widths, and no later biases has product of Frobenius norms
at most $P_0\ge1$. For every $0<\epsilon<1$, it has a same-depth
approximation $\widetilde z$ such that
\[
 \max_{x\in X}|z(x)-\widetilde z(x)|\le\epsilon,
 \qquad \prod_{j=1}^{D_0}\|\widetilde W^j\|_F\le2P_0,
\]
and every new hidden width is bounded by a universal polynomial in
\[
 P_0,D_0,R,\log(2|X|),\epsilon^{-1}.
\]
The bound is independent of the original widths and individual weights.
An augmented constant input is allowed and counts toward $R$.
\end{theorem}

\begin{proof}
Put $H=D_0-1$, $P_*=2P_0$, $e=\epsilon/(8H)$ and
$\eta=1/(4H)$. We compress the hidden layers in order from the first
to the last. At each step the product norm increases by at most
$1+\eta$, and the output changes by at most $6e$ on $X$.
Thus every intermediate product is below $P_*$, and the final product
is at most $e^{1/4}P_0<2P_0$, while total error is less than $\epsilon$.
A zero matrix makes the entire network zero and is trivial, so all
prefix and pair norms used below may be assumed positive. If a
compression step produces a zero matrix, stop and use the zero network
with width one; the accumulated approximation estimates still apply.

Consider a hidden matrix $W$ and its immediately following matrix $V$;
the latter may be the final linear readout. Let $S$ be the product of
all Frobenius norms strictly before $W$, with $S=1$ at the first layer,
and put
\[
 A=\|W\|_F\|V\|_F.
\]
The incoming feature vector $h(x)$ has norm at most $SR$. The part
strictly after $V$ has Lipschitz constant at most $P_*/(SA)$.
If $w_i^T$ is row $i$ of $W$ and $v_i$ is column $i$ of $V$, write
$a_i=\|w_i\|_2$ and $b_i=\|v_i\|_2$.
Zero rows or zero outgoing columns can be deleted.

Set
\[
 B=P_*/e,\qquad \tau^2=3e/(P_*R^3),\qquad
 \Delta=e/(P_*R),\qquad Q=(B/\tau)^2.
\]
Partition neurons into large, small and middle according as
$Sa_i>B$, $Sa_i<\tau$, or $\tau\le Sa_i\le B$.
There are at most $S^2\|W\|_F^2/B^2$ large neurons. Their activation
vector has norm at most the square root of this count. Deleting them
therefore changes the final output by at most $P_*/B=e$.

For the small neurons, replace tanh by the identity. Since
$|u-\tanh u|\le |u|^3/3$, their combined activation error has norm at
most $\tau^2R^3S\|W\|_F/3$. Propagation through the suffix bounds
the output error by $P_*\tau^2R^3/3=e$.
Their resulting contribution to the next preactivation is the linear map
\[
 B_0=\sum_{i\ \mathrm{small}}v_iw_i^T,
 \qquad \|B_0\|_*\le\sum_{i\ \mathrm{small}}a_ib_i,
\]
where $\|\cdot\|_*$ is the nuclear norm.

Let $A_m=\sum_{i\ \mathrm{middle}}a_ib_i$.
If $A_m<\Delta A$, delete all middle neurons: their output contribution
is at most $P_*R A_m/A\le e$. Otherwise $A_m\ge\Delta A>0$, and
we approximate their contribution by sampling. Define
\[
 Z=\sum_{i\ \mathrm{middle}}b_i/a_i,
 \qquad p_i=(b_i/a_i)/Z.
\]
Draw $m$ independent indices from $(p_i)$. A sampled neuron keeps its
incoming row $w_i^T$ and has outgoing column $v_i/(mp_i)$.
If $\widehat A=(Z/m)\sum_{r=1}^m a_{i_r}^2$ and $c=m/Z$, its new
incoming and outgoing squared Frobenius norms are exactly
\[
 c\widehat A\quad\hbox{and}\quad \widehat A/c.
\]
Moreover $\mathbb E\widehat A=A_m$. Since middle row norms lie between
$\tau/S$ and $B/S$, one has
\[
 Z\max_i a_i^2\le A_m Q\le AQ.
\]
The random vector $v_i\tanh(w_i^Th(x))/p_i$ has norm at most
$AQSR$, uniformly on $X$.

Consequently some choice of the samples has both
\[
 \widehat A\le A_m+\eta A
\]
and middle-preactivation approximation error at most $eSA/P_*$
simultaneously for every $x\in X$, provided
\[
 m\ge64Q^2\max\{\eta^{-2},(P_*R/e)^2\}
                       \log(8(|X|+1)).
\]
Here and below take integer ceilings. For detail, the scalar random
variables $Za_i^2/A$ are bounded by $Q$, so Hoeffding controls their
mean. For a Hilbert-valued random vector bounded by $AQSR$, the mean
deviation has expectation at most $AQSR/\sqrt m$; bounded differences
gives a dimension-free tail. A union bound over $X$ and the scalar
condition has positive success probability with the displayed constant.
Propagation of the vector error gives final output error at most $e$.

We next compress the linear branch. Truncate the singular-value
decomposition of $B_0$ to rank
\[
 k=\lceil P_*R/e\rceil,
\]
obtaining $B_k=UV_0$ with
$\|U\|_F^2=\|V_0\|_F^2=\nu:=\|B_k\|_*\le\|B_0\|_*$.
The residual operator norm is at most $\|B_0\|_* /(k+1)\le A/(k+1)$,
so its final output error is at most $P_*R/(k+1)\le e$.
If $\nu=0$, omit the linear branch entirely.

When the middle branch was sampled, retain $c=m/Z$. When it was
deleted and $\nu>0$, choose $c=1/(\nu S^2)$ instead.
Represent $B_k$ approximately using $K$ copies of each of its at most
$k$ singular channels: their incoming rows are the rows of
$\sqrt{c/K}\,V_0$, repeated $K$ times, and their outgoing columns are
the corresponding columns of $U/\sqrt{cK}$.
Their total incoming and outgoing squared norms are $c\nu$ and
$\nu/c$. Their linear limit is exactly $B_kh$.
The tanh linearization error at the next preactivation has norm at most
\[
 \frac{c\nu^2(SR)^3}{3K}.
\]
After propagation this is at most
$P_*c\nu S^2R^3/(3K)$, using $\nu\le A$.
In the sampled case,
\[
 Z\ge A_mS^2/B^2,
 \qquad c\nu S^2\le mB^2/\Delta;
\]
in the deleted case this quantity equals one. Hence it suffices to use
\[
 K=\left\lceil\frac{P_*R^3}{e}
                   \max\{1,mB^2/\Delta\}\right\rceil.
\]
The linear-channel approximation then contributes at most $e$ error.
When there is no sampled middle branch, the same bound is harmless;
one may use the displayed deterministic value of $m$ merely to define
a common width upper bound.

Combining the sampled and linear branches gives pair product norm
\[
 \widehat A+\nu
 \le A_m+\sum_{i\ \mathrm{small}}a_ib_i+\eta A
 \le(1+\eta)A.
\]
The last inequality is Cauchy--Schwarz over the original rows and
columns. If the middle branch was deleted, the pair product is simply
$\nu\le A$. These are exact norm calculations for the constructed
tanh matrices, not only bounds on their linear limits.

The compressed hidden width is at most $m+kK$. All quantities in this
bound are universal polynomials in $P_0,H,R,e^{-1}$ and
$\log(2|X|)$. Crucially, no original dimension of either adjacent hidden
layer occurs. Only the two adjacent matrices are altered. Processing
the next hidden layer changes its own width without undoing the width
bound for the preceding one. Uniform zero-padding at the end makes
all hidden widths equal if desired. Summing the at most six errors and
multiplying the pair-product bounds proves the theorem.
\end{proof}

\begin{corollary}[Uniform compression on a bounded input set]
\label{cor:pf-compact-compression}
If the input dimension is $q$ and $X$ is any subset of the Euclidean
ball of radius $R\ge1$, the same conclusion holds with width polynomial
in $q,P_0,D_0,R,\epsilon^{-1}$, independently of whether $X$ is finite.
\end{corollary}

\begin{proof}
Choose an $\epsilon/(6P_0)$-net of the radius-$R$ ball with at most
$(1+12P_0R/\epsilon)^q$ points. Apply the finite-set theorem to this net
with error $\epsilon/2$. The original and compressed networks have
Lipschitz constants at most $P_0$ and $2P_0$, respectively. Their
difference therefore changes by at most $\epsilon/2$ between any point
and its nearest net point. This proves uniform error at most $\epsilon$.
The logarithm of the net size is polynomial in the displayed quantities.
\end{proof}

\begin{corollary}[A width-free sufficient representation condition]
\label{cor:width-free-pf-sufficiency}
Suppose $f$ has some finite-width, depth-$D_0$ tanh representation with
unit global margin and product of Frobenius norms at most $P_0\ge1$.
For each fixed $\zeta>0$, there are universal polynomials and a universal
constant $c$ such that its unrestricted Gaussian posterior sample
complexity is polynomial in
\[
 d,D_0,P_0,L,\varepsilon^{-1},\log(2/\delta)
\]
whenever
\[
 L\ge D_0+c\log(2+d+D_0+P_0),
\]
simultaneously for every $N\ge C_L(f)$ and every
$0<\sigma^2\le N^{-(1+\zeta)}$.
The same conclusion holds at every $L\ge2D_0$.
There is no dependence on the original width or on individual weights.
\end{corollary}

\begin{proof}
Apply Theorem~\ref{thm:pf-width-compression} with the augmented Boolean
cube, $R=\sqrt{d+1}$, $|X|=2^d$, and approximation error $1/4$.
The resulting representation has width polynomial in $d,D_0,P_0$,
the same depth, and global signed margin at least $3/4$.
Double its readout to obtain unit margin and product norm at most
$4P_0$. Theorem~\ref{thm:product-target-log-padding} then applies.
Its $T$ is polynomial in $d,D_0,P_0$, so the required logarithmic
overhead has the displayed form after increasing a universal constant.
\end{proof}

Together with the posterior ceiling, this identifies polynomial product
norm as a substantive representation criterion for this prior. The
necessity statement concerns constant-error approximations at the same
learner depth; the sufficiency statement above assumes an exact Boolean
representation and permits logarithmically many added layers. These
qualifications prevent treating the two statements as an exact
equivalence.
